Вы должны послать ровно один файл, названный \texttt{friend.cpp}. В этом файле должна быть реализована функция, описанная выше с указанными ниже прототипами. На языках C/C++ вы должны подключить заголовочный файл \texttt{friend.h}.

\begin{center}
\renewcommand{\arraystretch}{1.5}
\begin{tabular}{|c|c|c|c|}
\hline
Этап & Хозяин & Протокол & Добавленные пары друзей \\
\hline
1 &  0 & IAmYourFriend &  (1, 0) \\
\hline
2 & 0 & MyFriendsAreYourFriends & (2, 1) \\
\hline
3 & 1 & WeAreYourFriends & (3, 1), (3, 0), (3, 2)  \\
\hline
4 & 2 & MyFriendsAreYourFriends & (4, 1), (4, 3) \\
\hline
5 & 0 & IAmYourFriend & (5, 0) \\
\hline
\end{tabular}
\end{center}

Вначале сеть содержит только участника с номером 0. Хозяин первого этапа (участник с
номером 0) приглашает нового участника с номером 1, используя протокол \texttt{IAmYourFriend}, и они становятся друзьями. Хозяин второго этапа (снова участник с номером 0) приглашает второго участника, используя протокол \texttt{MyFriendsAreYourFriends}, который делает участника с номером 1 (единственный друг хозяина) единственным другом участника с номером 2. Хозяин третьего этапа (участник с номером 1) добавляет третьего участника, используя протокол \texttt{WeAreYourFriends}, что делает третьего участника другом первого участника (хозяина) и участников с номерами 0 и 2 (они же друзья хозяина). Этапы 4 и 5 также показаны в таблице выше. Финальная сеть представлена на рисунке ниже, где число в кружке показывает номер участника, а число рядом с кружком показывает уровень доверия в опросах для этого участника. Выборка, состоящая из участников с номерами 3 и 5, имеет суммарный уровень доверия в опросах, который составляет $20 + 15 = 35$, что является максимально возможным
суммарным уровнем доверия.

\includegraphics{1.png}